\begin{ex}
Góc nhọn là góc bé hơn góc vuông đúng hay sai? (Đúng ghi 1, sai ghi 0).
\shortans{1}
\loigiai{
Đúng, góc nhọn bé hơn góc vuông.
}
\end{ex}

\begin{ex}
Góc tù là góc lớn hơn góc vuông và bé hơn góc bẹt đúng hay sai? (Đúng ghi 1, sai ghi 0).
\shortans{1}
\loigiai{
Đúng.
}
\end{ex}

\begin{ex}
Góc bẹt bằng mấy góc vuông?
\shortans{2}
\loigiai{
Góc bẹt bằng 2 góc vuông ($180^\circ = 2 \times 90^\circ$).
}
\end{ex}

\begin{ex}
Một hình tam giác có thể có nhiều nhất bao nhiêu góc tù?
\shortans{1}
\loigiai{
Tam giác chỉ có nhiều nhất 1 góc tù.
}
\end{ex}

\begin{ex}
Một hình tam giác có thể có nhiều nhất bao nhiêu góc vuông?
\shortans{1}
\loigiai{
Tam giác chỉ có nhiều nhất 1 góc vuông.
}
\end{ex}

\begin{ex}
Một hình tam giác nhọn có bao nhiêu góc nhọn?
\shortans{3}
\loigiai{
Có 3 góc nhọn.
}
\end{ex}

\begin{ex}
Hình chữ nhật có bao nhiêu góc vuông?
\shortans{4}
\loigiai{
Có 4 góc vuông.
}
\end{ex}

\begin{ex}
Hình vuông có bao nhiêu góc vuông?
\shortans{4}
\loigiai{
Có 4 góc vuông.
}
\end{ex}

\begin{ex}
Góc tạo bởi hai kim đồng hồ lúc 2 giờ đúng là góc gì? (Góc nhọn ghi 1, góc tù ghi 2, góc bẹt ghi 3).
\shortans{1}
\loigiai{
Góc $60^\circ$ là góc nhọn (1).
}
\end{ex}

\begin{ex}
Góc tạo bởi hai kim đồng hồ lúc 4 giờ đúng là góc gì? (Góc nhọn ghi 1, góc tù ghi 2, góc bẹt ghi 3).
\shortans{2}
\loigiai{
Góc $120^\circ$ là góc tù (2).
}
\end{ex}

\begin{ex}
Góc tạo bởi hai kim đồng hồ lúc 6 giờ đúng là góc gì? (Góc nhọn ghi 1, góc tù ghi 2, góc bẹt ghi 3).
\shortans{3}
\loigiai{
Góc $180^\circ$ là góc bẹt (3).
}
\end{ex}

\begin{ex}
Trong các góc sau, góc nào là góc tù: $45^\circ, 90^\circ, 120^\circ, 180^\circ$? (Điền số đo độ của góc đó).
\shortans{120}
\loigiai{
Góc tù là $120^\circ$.
}
\end{ex}

\begin{ex}
Trong các góc sau, góc nào là góc nhọn: $30^\circ, 90^\circ, 105^\circ, 180^\circ$? (Điền số đo độ của góc đó).
\shortans{30}
\loigiai{
Góc nhọn là $30^\circ$.
}
\end{ex}

\begin{ex}
Một chiếc quạt xòe tròn một nửa tạo thành hình quạt với góc ở tâm là góc gì? (Góc bẹt ghi 1, vuông ghi 2).
\shortans{1}
\loigiai{
Đó là góc bẹt ($180^\circ$).
}
\end{ex}

\begin{ex}
Cho một góc vuông, vẽ một tia nằm giữa hai cạnh của góc vuông. Ta được mấy góc nhọn?
\shortans{2}
\loigiai{
Tia phân chia góc vuông thành 2 góc nhọn.
}
\end{ex}

\begin{ex}
Một hình thoi có 2 góc nhọn và mấy góc tù?
\shortans{2}
\loigiai{
Hình thoi có 2 góc nhọn và 2 góc tù.
}
\end{ex}

\begin{ex}
Một hình bình hành có 2 góc nhọn và mấy góc tù?
\shortans{2}
\loigiai{
Có 2 góc tù.
}
\end{ex}

\begin{ex}
Góc $89^\circ$ là góc gì? (Góc nhọn ghi 1, góc tù ghi 2).
\shortans{1}
\loigiai{
Góc $89^\circ < 90^\circ$ là góc nhọn.
}
\end{ex}

\begin{ex}
Góc $91^\circ$ là góc gì? (Góc nhọn ghi 1, góc tù ghi 2).
\shortans{2}
\loigiai{
Góc $91^\circ > 90^\circ$ là góc tù.
}
\end{ex}

\begin{ex}
Góc đỉnh O có hai cạnh OA và OB thẳng hàng ngược hướng nhau là góc gì? (Góc bẹt ghi 1, góc vuông ghi 2).
\shortans{1}
\loigiai{
Đó là góc bẹt.
}
\end{ex}

\begin{ex}
Có bao nhiêu góc nhọn trong một hình tam giác vuông?
\shortans{2}
\loigiai{
Tam giác vuông có 1 góc vuông và 2 góc nhọn.
}
\end{ex}

\begin{ex}
Có bao nhiêu góc nhọn trong một hình tam giác tù?
\shortans{2}
\loigiai{
Tam giác tù có 1 góc tù và 2 góc nhọn.
}
\end{ex}

\begin{ex}
Lúc 5 giờ đúng, góc tạo bởi kim giờ và kim phút là góc gì? (Góc nhọn ghi 1, góc tù ghi 2).
\shortans{2}
\loigiai{
Lúc 5 giờ góc bằng $150^\circ$ là góc tù (2).
}
\end{ex}

\begin{ex}
Lúc 1 giờ đúng, góc tạo bởi kim giờ và kim phút là góc gì? (Góc nhọn ghi 1, góc tù ghi 2).
\shortans{1}
\loigiai{
Lúc 1 giờ góc bằng $30^\circ$ là góc nhọn (1).
}
\end{ex}

\begin{ex}
Góc bẹt có số đo bằng bao nhiêu độ?
\shortans{180}
\loigiai{
Góc bẹt bằng $180^\circ$.
}
\end{ex}
